% ---------------------------------------------------------------------------------------------------------
% 評価の規約（実験 0）
% ---------------------------------------------------------------------------------------------------------
\section{評価の規約}
\label{sec:protocol}
閉ループの実験に入る前に，環境，装置，接触モデル，体格，成功の定義を固定した．以降の実験はすべてこの規約に従う．

\subsection{環境}
第~\ref{sec:model}節の縮約モデルの物理（A に拘束された振り，自由飛行，捕捉，摩擦・引っ掛かり錐，可動域，壁）を
CasADi からコード生成してコンパイルした環境を用いる．最適化器と同じ記号力学を共有するので，
拘束された振りの加速度は最適化器の力学と相対誤差 $10^{-3}$ で一致する．
制御周期は 2\,ms，積分は 0.5\,ms 刻みの 4 次 Runge--Kutta 法である．
関節の可動域は軟らかいストップ（剛性 $15\,\taucap_j$/rad）で表し，指令トルクには最適化器と同じ変化率の上限を課す．
捕捉は，掛かり線が B の先端の手前 2.5\,cm から壁までの範囲に許容速度で入ったとき，指が反射的に閉じて有限剛性のフックになるものとした．
フックに求められる力が錐を出ると手は突起の上を滑り，先端を越えると外れる．
指のフックは，摩擦・引っ掛かり錐に加えて，壁から離れる向きに $0.1\Fref$ の機械的な引っ掛かりをもつ．
B に掛かった後は，体が B の壁板に触れて突っ張ることを許す（ばね・ダンパの接触，10\,cm を超える食い込みだけを失敗とする）．
振りと飛行の間に体が壁に 2\,cm を超えて食い込んだら失敗とする．

離手した後の飛行と捕捉は，全手法に共通の\textbf{着地反射}が担う．
これは，保存してある参照解のうち体格の最も近いものの飛行中の関節角を PD で追従し，手先を B の先端へ寄せる補正を加え，
捕捉後は参照解の保持の関節角を追従する制御器である．
これにより，比較の対象を「振りで勢いをつけることと，いつ離手するか」に限定する．

\subsection{成功の定義と失敗の分類}
主指標は次の条件をすべて満たした試行の割合である：
A にぶら下がって静止した状態から始め，期限（1 周期 $T$）までに離手し，B を捕捉し，2\,s 保持する．
その間ずっと $\|\bm R(t)\|\le\Fmax$ であり，関節と接触の制約を破らない．
$\Fmax=c\,\Fref$ での成功率を $S_c$ と書き，$c\in\{1.0,1.25,1.5,1.75,2.0,2.25,2.5\}$ について報告する．
見出しの値は $S_{2.0}$（$\Fmax=2600$\,N）である．
試行は $\Fmax=2.5\Fref$ で実行し，各水準を初めて超えた時刻を記録しておく．
ある水準での成否は「その水準を超える前に成功の条件を満たしたか」で決まるので，低い水準の成功率は記録から正確に求まる．

失敗した試行は，最初に起きた失敗の種類と時刻で分類する：
A から滑落，壁に接触，離手せず（期限切れ），B を掴み損ね，B の前面に衝突，保持中に指が外れる，把持容量超過．
振りの途中の状態から始める試行は，主指標には含めず，\textbf{回復試験}として別に報告する．

開始状態の摂動は，静止姿勢の体幹の角度と 4 つの関節角に標準偏差 $\sigma_\theta$ の，その角速度に $\sigma_{\dot\theta}$ の正規乱数を加えて作る
（壁に食い込む姿勢は引き直す）．「弱い摂動」は $\sigma_\theta=0.01$\,rad，$\sigma_{\dot\theta}=0.1$\,rad/s，
「強い摂動」は $\sigma_\theta=0.03$\,rad，$\sigma_{\dot\theta}=0.2$\,rad/s である．

\subsection{比較する制御器}
比較する制御器を表~\ref{tab:methods} に示す．
中心となる比較は，\textbf{双対場 MPC と価値 MPC の対}である．
両者は MPC の定式化，学習データ，ネットワークの構造，離手のタイマー，把持容量の上限が同じで，
終端の価値を学習するときに costate ラベルを使ったかどうかだけが違う．
成功率の差 $\Delta S=S_{2.0}(\text{双対場 MPC})-S_{2.0}(\text{価値 MPC})$ を，同じ体格・同じ摂動・同じ学習の種で対にした試行について求め，
ブートストラップによる 95\,\%信頼区間を付ける．
オラクル MPC+ は自分の計画の飛行と保持を着地反射に渡すので，共通の着地反射を使うという条件を満たさない．
したがって比較の対象ではなく，到達できる水準の参考として示す．
\begin{table*}[t]
\centering\scriptsize
\caption{比較する制御器．すべて同じ環境，同じ $\Fmax$，同じ実行系（制御周期 2\,ms）で評価する．着地反射の「共通」は保存した参照解を使うことを表す．}
\label{tab:methods}
\resizebox{\textwidth}{!}{%
\begin{tabular}{llp{0.56\textwidth}ll}
\toprule
名称 & 略称 & 振りの制御 & 離手の決め方 & 着地反射 \\
\midrule
再生 & REPLAY & 参照解のトルクをそのまま加え，参照の関節角を弱い PD（1/rad，0.1\,s/rad）で追従する & 参照解の時刻 & 共通 \\
行動模倣 & BC & オラクルのトルクを状態から回帰したネットワーク & 残り時間の予測 & 共通 \\
追従 MPC & TRACK & 式~\eqref{eq:mpc} と同じ MPC で，終端の価値の代わりに参照解の状態への追従誤差を最小化する & 参照解の時刻 & 共通 \\
価値 MPC & VMPC & 式~\eqref{eq:mpc}．終端は costate ラベルなしで学習した価値（$\alpha=0$） & 残り時間の予測 & 共通 \\
双対場 MPC & SMPC & 式~\eqref{eq:mpc}．終端は Sobolev 学習した価値（$\alpha=1$） & 残り時間の予測 & 共通 \\
オラクル MPC & ORACLE & 残り最適化問題を 0.5\,s ごと（と離手の 0.25\,s 前）に解き直し，その計画を追従する & 計画の時刻 & 共通 \\
オラクル MPC+ & ORACLE+ & 同上 & 計画の時刻 & 自分の計画（参考） \\
\bottomrule
\end{tabular}}
\end{table*}
